\documentclass[preprint,12pt]{elsarticle}

\usepackage{amsmath}
\usepackage{amssymb}
\usepackage{booktabs}
\usepackage{enumitem}
\usepackage{graphicx}
\usepackage[hyphens]{url}
\usepackage{hyperref}
\usepackage{makecell}
\usepackage{microtype}
\usepackage{tabularx}
\usepackage{xcolor}
\usepackage{multirow}

\journal{Computers \& Security}

\begin{document}

\begin{frontmatter}

\title{Toward Evidence-Grounded Skill Evolution for\\
LLM-Based Vulnerability Analysis}

\author[inst1]{Author One}
\ead{author1@example.com}
\author[inst1]{Author Two}
\ead{author2@example.com}
\address[inst1]{Affiliation, City, Country}

\begin{abstract}
Reusable natural-language skills are becoming a practical control surface
for LLM agents, but in vulnerability analysis the key problem is not merely
how to generate candidate skills.  The harder problem is deciding when
evidence from a completed analysis run is strong enough to justify changing
the live skill library that future analyses will depend on.  We study this
problem as \emph{evidence-constrained skill evolution} for
vulnerability-analysis agents.  Our setting couples blind remote analysis
with post-run external confirmation: an agent solves a task on an isolated
virtual machine, the local system scores the final answer and checks it
against source predicates, logic harnesses, behavioral oracles, or
sanitizer-backed evidence, and only then may an evolution service propose
and validate candidate skill updates.  We present a split-plane
architecture that separates analysis from confirmation, a run archival
scheme that links selected skills to immutable session segments, a
lightweight run-level feedback interface for candidate generation, and a
publication gate that stages candidate skills before they can enter the
live library.  Empirically, our primary 36-run blind A/B study across six
firmware cases shows strongly heterogeneous skill effects, with mean score
changes ranging from $+1.67$ to $-1.00$.  Complementary targeted firmware
ablations on two CGI cases further show that skill structure matters: a
branch-first seed skill is more stable than a broader but less focused live
profile, while wrong or absent skills can still yield deceptively plausible
answers.  Across 94 archived gate decisions, candidate generation is not
the main bottleneck; admission into the live library is.  Replay-based
gates can reject clearly poor candidates, but they still struggle to
separate genuine improvement from harmless redundancy when the baseline
model already performs well.  The resulting contribution is a
vulnerability-analysis-oriented research scaffold for studying when
validated runs should, and should not, become reusable skills.
\end{abstract}

\begin{keyword}
LLM agent \sep skill evolution \sep vulnerability analysis \sep
external validation \sep blind benchmark \sep evidence-grounded feedback
\end{keyword}

\end{frontmatter}

\section{Introduction}
\label{sec:introduction}

LLM agents increasingly use external memories, reusable instructions, and
natural-language skills to improve over repeated tasks.  Systems such as
Voyager~\cite{voyager} and ExpeL~\cite{expel} show that agents can
accumulate procedural knowledge from interaction and reuse it later.
SkillClaw~\cite{skillclaw} extends this direction to a shared skill
ecosystem, while recent work studies context-to-skill induction
(\cite{ctx2skill}), on-policy skill distillation (\cite{opid}),
environment-grounded skill formation (\cite{spark,skillevolbench}), and
verifier-guided skill refinement (\cite{trace2skill,veriskill}).

At the same time, two recent findings complicate the optimistic picture
that ``more skills'' naturally means ``better agents.''  First,
skill-induced failures are common: even topically relevant skills can push
an agent onto the wrong procedure or the wrong abstraction
level~\cite{skillharmful}.  Second, self-improving agents can accumulate
unsafe or low-value procedural knowledge if successful runs are converted
into persistent skills without enough scrutiny~\cite{misevolution}.  In
parallel, security-specific work has begun to evolve reusable playbooks for
auditing~\cite{evohunt} and to construct more realistic long-horizon
security benchmarks~\cite{secbenchpro}.  These lines of work establish that
skill evolution, agentic security, and realistic evaluation are all active
research directions.  What remains less clear is how they should be joined
in a deployment-oriented vulnerability-analysis loop.

That integration problem is the starting point of this paper.  In a real
vulnerability-analysis workflow, the difficult question is not whether an
LLM or agent can produce a candidate skill revision.  The difficult
question is whether a completed run contains reusable procedural knowledge
that deserves to enter the \emph{live} skill library seen by later runs.
This question is more stringent than ordinary task success.  A run can be
correct yet redundant, correct yet overfit to a single instance, or
correct yet unsafe to generalize.  The engineering surface therefore
includes not only skill retrieval and revision, but also run archival,
evidence collection, candidate staging, and admission control.

Vulnerability analysis is a particularly instructive domain for studying
this problem because it provides something many generic agent tasks do not:
\emph{external, executable confirmation}.  Depending on the case, a claim
can be checked by source predicates, logic harnesses, behavioral oracles,
or sanitizer-backed crashes.  These signals do not prove causal skill
contribution, but they do let us constrain the evolution loop using
evidence that is independent of the model's own narrative.  This suggests a
different research focus from prior work on skill generation or pure
retrieval: rather than asking only how to produce better skills, we can ask
how to govern the transition from a \emph{validated run} to a \emph{live
reusable skill}.

Our goal is therefore narrower and more operational than a general theory
of skill learning.  We do not claim a new retrieval algorithm, a new
revision model, or a complete solution to causal attribution.  Instead, we
study a full vulnerability-analysis loop in which a blind run executes on a
remote analysis environment, selected skills are injected server-side, the
final answer is scored and externally confirmed on a separate local plane,
the resulting record is transformed into a candidate update, and the update
is published only if it passes an explicit gate.  This framing lets us ask
three concrete questions: how unstable skill effects are in blind security
tasks; whether lightweight post-run feedback is useful enough to drive
candidate generation; and whether publication control, rather than
candidate generation, is the real bottleneck in practice.

This paper makes four contributions:

\begin{enumerate}[leftmargin=*]
    \item \textbf{An evidence-constrained lifecycle view of skill evolution
    for vulnerability-analysis agents.} We formulate the transition from a
    blind completed run to a reusable live skill as a structured lifecycle
    involving execution, confirmation, archival, candidate generation, and
    publication control.

    \item \textbf{A split-plane architecture for blind analysis and local
    confirmation.} We design a remote-analysis / local-confirmation
    separation that preserves blind execution while still supporting
    post-run evidence processing and evolution.

    \item \textbf{A problem formulation and empirical analysis of live skill
    publication control.} We show that the main bottleneck in the current
    loop is not candidate generation but deciding whether a candidate skill
    should enter the live library.

    \item \textbf{An experimental protocol for measuring heterogeneous skill
    effects and skill necessity tendencies in vulnerability analysis.}
    Across a 36-run blind A/B study over six firmware cases and targeted
    profile ablations on two CGI tasks, we show that skill effects are
    strongly case-dependent and sensitive to skill structure rather than
    topic match alone.
\end{enumerate}

The intended scope of this paper is therefore a \emph{framework paper with
initial empirical evidence}.  Our objective is to establish the research
problem, the system structure, and the first round of measurable findings.
We do not claim that autonomous skill improvement has already been solved.
Instead, we argue that the field needs a cleaner way to study when
validated security runs should, and should not, become persistent skills,
and we present a first operational scaffold for that study.

The remainder of the paper is organized as follows.
Section~\ref{sec:problem} formalizes the skill-evolution problem and its
two key subproblems: run-level attribution and safe publication.
Section~\ref{sec:method} presents the framework.
Section~\ref{sec:implementation} describes the prototype.
Section~\ref{sec:evaluation} reports experimental methodology and results.
Section~\ref{sec:discussion} discusses implications.
Section~\ref{sec:related} surveys related work.
Section~\ref{sec:threats} addresses threats to validity.
Section~\ref{sec:conclusion} concludes.

\section{Problem Formulation}
\label{sec:problem}

\subsection{Task Run and Validation Vector}

Let a task run be
\begin{equation}
    r = (x, S_r, \tau, y, E),
\end{equation}
where $x$ is a blind vulnerability-analysis task, $S_r$ is the set of
skills selected and injected for the run, $\tau$ is the agent trajectory,
$y$ is the final answer (predicted files, functions, root cause, and
evidence), and $E$ is externally collected evidence.

We retain a validation vector rather than a single success label:
\begin{equation}
    V(r) = \left(q_{\mathrm{task}},\; q_{\mathrm{confirm}},\;
    a_{r,s}\right).
\end{equation}
Here $q_{\mathrm{task}}$ measures whether the answer localizes and explains
the ground-truth vulnerability; $q_{\mathrm{confirm}}$ represents the
strength of external source or runtime evidence; and $a_{r,s}$ records
\emph{observed relevance} between skill $s$ and the task.

We deliberately call $a_{r,s}$ \emph{observed relevance} rather than
\emph{causal contribution}.  This distinction is the central problem of
this paper.

\subsection{The Skill Attribution Problem}
\label{sec:attribution_problem}

Suppose a run achieves high $q_{\mathrm{task}}$ and high
$q_{\mathrm{confirm}}$: the agent correctly identified the vulnerability
file, function, and root cause, and the result was confirmed by a sanitizer
crash or source predicate.  Can we conclude that the injected skill $s$
causally contributed to this success?

In general, no.  Consider the following decomposition.  Let $y_0$ be the
answer the model would produce \emph{without} skill $s$ (the counterfactual
baseline), and let $y_s$ be the answer \emph{with} skill $s$.  The causal
effect of $s$ on the task outcome is:
\begin{equation}
    \Delta(s) = q_{\mathrm{task}}(y_s) - q_{\mathrm{task}}(y_0).
    \label{eq:causal_effect}
\end{equation}
If $\Delta(s) > 0$, the skill helped.  If $\Delta(s) \leq 0$, it did not
help or was harmful.  However, computing $\Delta(s)$ requires running the
agent twice---once with $s$ and once without---which doubles the
computational cost per attribution.  More importantly, LLM outputs are
stochastic: a single pair of runs may not reflect the true expected
effect, requiring multiple repeated trials for statistical power.

Prior work addresses this through counterfactual ablation.  Differential
comparison between a skill-guided run and a no-skill reference
run~\cite{skillharmful}, and Shapley-style step valuation over skill
traces~\cite{skillshapley,shapley}, both provide stronger attribution
signals than single-run observation.  However, Shapley-style estimation
requires evaluating many subsets of skill steps, making it expensive for
long skills.

The key question we investigate is: \emph{can executable domain evidence
provide a cheaper approximation of $\Delta(s)$ without running
counterfactual ablations?}

\subsection{The Safe Skill Publication Problem}
\label{sec:publication_problem}

When a run is completed and validated, the system may generate a candidate
skill revision $\Delta s$ derived from the run's trajectory and feedback.
The question is whether to publish $\Delta s$ into the live skill library
that subsequent agents will use.

Following the principle that ``evolution optimizes task outcomes rather
than procedure safety''~\cite{misevolution}, we require a publication gate:
\begin{equation}
    \mathrm{Publish}(\Delta s) =
    \mathbb{I}\left[
      \mathrm{EvidenceSufficient}(\Delta s)
      \;\land\;
      \mathrm{FollowupPass}(\Delta s)
    \right].
    \label{eq:gate}
\end{equation}
Here $\mathrm{EvidenceSufficient}$ checks whether the originating run had
adequate evidence to justify a skill update, and $\mathrm{FollowupPass}$
runs a follow-up validation with the candidate skill to check whether it
maintains or improves performance.

The structural challenge is in $\mathrm{FollowupPass}$.  In a
\emph{replay-only} paradigm, the follow-up runs the candidate skill on the
same or similar task and checks whether the result passes validation.  But
this only measures $q_{\mathrm{task}}(y_{\Delta s})$, not
$\Delta(\Delta s) = q_{\mathrm{task}}(y_{\Delta s}) -
q_{\mathrm{task}}(y_0)$.  Without the counterfactual baseline $y_0$, the
gate cannot distinguish ``the skill helped'' from ``the model would have
succeeded anyway.''  When the gate requires \emph{strict improvement}
(mean score with candidate $>$ mean score without candidate), this leads to
systematic rejection of candidates that are not harmful but merely
redundant---a high false-negative rate.

\subsection{Research Hypotheses}

The study is guided by three hypotheses:

\begin{description}
    \item[H1: Skill instability.] Injected domain skills produce unstable
    effects on vulnerability-analysis quality, including net harm in some
    cases.
    \item[H2: Evidence approximation.] Executable domain evidence can
    approximate skill attribution at lower cost than counterfactual
    ablation, but with reduced precision.
    \item[H3: Gate limitation.] A replay-only publication gate suffers a
    structural false-negative limitation that cannot be resolved by
    parameter tuning alone.
\end{description}

\section{Evidence-Constrained Skill Evolution Framework}
\label{sec:method}

We propose a framework that connects blind task execution, skill injection
auditing, external evidence confirmation, evidence-conditioned feedback,
and gated candidate publication.  The framework is designed around a simple
principle: in vulnerability analysis, post-run evidence should constrain
both what we infer about a skill and what we allow the system to retain for
future use.  The goal is not to make the evolution service omniscient.  The
goal is to ensure that the path from a completed run to a persistent live
skill is explicit, inspectable, and conservative enough for security work.

\subsection{Leakage-Controlled Blind Execution}
\label{sec:blind}

Each evaluation begins from a benchmark specification containing private
ground truth and a separately rendered blind prompt.  The agent receives
the target workspace and the analysis objective, but not the answer fields,
confirmation scripts, or solution-bearing artifact names.  The benchmark
controller retains the expected file, function, root cause, and
confirmation policy.

The analysis runs on a remote virtual machine, separated from the
benchmark ground truth and from the local service that evaluates the
result.  This separation is methodologically important: it prevents
inadvertent answer leakage through shared filesystems or local environment
artifacts, and it mirrors real-world deployment where the analyst's
environment is isolated from the vulnerability database.

\subsection{Auditable Skill Injection}
\label{sec:injection}

For each run, the system records the selected skill names, the injected
prompt content hash, the request session identifier, and the immutable
session segment consumed by the evolution process.

This recording addresses a prerequisite for studying skill evolution.
Any claim about a skill's effect is uninterpretable if the intended skill
cannot be linked to a specific request segment, or if the injected content
cannot be reconstructed after the run.  Injection auditing therefore does
not solve attribution, but it narrows the question from ``was the skill
present?'' to the more meaningful ``what procedural prior was exposed to
the model in this specific run, and how should that exposure be archived
for later judgment?''

\subsection{Multi-Evidence External Confirmation}
\label{sec:confirmation}

The agent's output is evaluated along two complementary dimensions.  First,
\emph{structured scoring} checks the predicted vulnerability identity,
file, function, root cause, and cited evidence against ground truth.
Second, \emph{case validators} test claims against the target environment.

We classify evidence into four categories, each making a different claim
about the vulnerability:

\begin{itemize}[leftmargin=*]
    \item \textbf{Source-backed evidence:} verifies a concrete code
    predicate or boundary relation (e.g., an unsafe function call exists
    at the predicted location).
    \item \textbf{Behavior-backed evidence:} reproduces an expected program
    path or output condition (e.g., the binary processes the input and
    reaches the predicted function).
    \item \textbf{Logic-backed evidence:} exercises a controlled harness
    for the vulnerable state transition (e.g., a test input triggers the
    expected code path without a full crash).
    \item \textbf{Sanitizer-backed evidence:} reproduces a memory-safety
    failure with relevant runtime diagnostics (e.g., AddressSanitizer
    reports a buffer overflow at the predicted location).
\end{itemize}

The distinction matters because different evidence classes support
different claims about the analysis.  A sanitizer crash is stronger
evidence for a memory-safety vulnerability than a source-code predicate
alone.  The framework records the accepted evidence class rather than
mapping every passing script to the same notion of ``confirmed.''

\subsection{Evidence-Conditioned Feedback}
\label{sec:feedback}

After validation, the system constructs feedback at the run--skill level.
This feedback is an operational interface for evolution rather than a
standalone causal estimator.  Its role is to summarize the relation between
task outcome, external confirmation, and observed skill relevance in a form
that can drive candidate generation and later gate decisions.  In that
sense, it is a low-cost approximation to the latent quantity $\Delta(s)$ in
Equation~\ref{eq:causal_effect}, but we do not present it as a substitute
for full counterfactual attribution.

A skill receives one of four statuses:

\begin{description}[style=nextline,leftmargin=2.4cm]
    \item[Supported ($\hat{\Delta} > 0$):] the task is externally validated
    and the skill is aligned with the analysis procedure (i.e., the skill's
    described steps match the agent's actual trajectory).  This is
    positive evidence for the skill, but not proof of causality.
    \item[Mismatched ($\hat{\Delta} \approx 0$):] the task may be correct,
    but the skill is unrelated to the target or method.  The model likely
    succeeded on its own; the skill was redundant.
    \item[Contradicted ($\hat{\Delta} < 0$):] the skill appears to drive a
    false localization, unsupported claim, or failed confirmation.  The
    skill was actively harmful.
    \item[Unknown:] the available trace cannot distinguish contribution
    from coincidence.
\end{description}

The key idea is practical rather than metaphysical.  Vulnerability-analysis
workflows already produce task scores, evidence classes, and auditable
records of the selected skill and final answer.  We reuse these signals to
derive an attribution-oriented label at the cost of one run, whereas
counterfactual ablation requires at least two runs and Shapley-style
valuation~\cite{skillshapley} can be exponentially expensive in the number
of skill steps.  This makes the label useful as a triage signal for
evolution even when it remains imperfect as scientific evidence of causal
contribution.

However, this approximation has lower precision.  A ``supported'' skill may
be supported only by coincidence: the model might have succeeded regardless.
A ``mismatched'' skill might have subtly helped in ways the alignment check
cannot detect.  We empirically evaluate this precision trade-off in
Section~\ref{sec:results}.

\subsection{Candidate-Level Publication Gate}
\label{sec:gate}

Evolution produces a candidate revision $\Delta s$ rather than directly
overwriting the live skill.  Publication is gated by
Equation~\ref{eq:gate}.  In our prototype, follow-up validation replays the
candidate on the same or a closely related task and checks whether the
candidate remains consistent with the stored evidence.

The most conservative form of this gate is a \emph{strict improvement}
policy: the candidate must achieve a mean validation score that exceeds a
baseline score by a margin.  This design is motivated by the concern that
successful-but-unsafe or merely coincidental experiences should not be
persisted as reusable skills.

However, as formalized in Section~\ref{sec:publication_problem}, a
replay-only gate has a structural limitation.  The follow-up measures
$q_{\mathrm{task}}(y_{\Delta s})$ but not
$\Delta(\Delta s) = q_{\mathrm{task}}(y_{\Delta s}) -
q_{\mathrm{task}}(y_0)$.  When the baseline is high (the model can already
solve the task), strict improvement is nearly impossible to establish
because the candidate's replay score is bounded by the same ceiling.  We
provide empirical evidence of this limitation in
Section~\ref{sec:results_gate}.

\subsection{Design Rationale: Live, Candidate, and Share Separation}

Live skills, candidate revisions, and shared runtime state are stored in
separate directories.  This separation is not merely a
software-organization choice; it is part of the experimental design.  It
ensures that candidate skills can be generated, inspected, replayed, and
rejected without silently changing the library seen by the next run.

\section{Prototype Realization}
\label{sec:implementation}

We realize the framework as a research extension to
SkillClaw~\cite{skillclaw}.  The original system already provides request
proxying, skill retrieval and injection, session capture, skill sharing,
and an evolution service.  Our extension does not replace these core
services with a new retrieval or revision algorithm.  Instead, it adds the
elements needed to study the run-to-library transition in a security
setting: leakage-controlled benchmark cases, remote blind execution,
heterogeneous vulnerability validators, candidate staging, and a structured
handoff from finalized runs to validation-gated evolution.

\subsection{Analysis Plane and Evidence Plane}

The runtime is divided into two planes.  The \emph{analysis plane} consists
of an LLM agent client on a remote VM that analyzes the target through the
SkillClaw proxy.  The \emph{evidence plane} is a local service that retains
private ground truth, finalizes the run record, executes validators, and
sends a closed session/feedback pair to the evolution service.  This split
preserves the blind-evaluation setting while still allowing rich local
post-processing and confirmation.

Client-provided session identifiers are preserved when available, and long
sessions are archived as incrementally consumable immutable segments.  This
ensures that later interaction can extend a session without rewriting
feedback already consumed by the evolver.

\subsection{Benchmark Cases}

We define benchmark cases in JSON specifications containing: the blind
workspace path, the analysis objective (rendered without CVE identifiers or
solution-bearing names), private ground truth (expected files, functions,
root cause, required evidence terms), and confirmation policies.

The current benchmark includes 13 cases spanning two categories:
\begin{itemize}[leftmargin=*]
    \item \textbf{Firmware cases (8):} embedded HTTPd command injection
    and buffer overflow vulnerabilities in router firmware binaries from
    multiple vendors, analyzed via static reverse engineering.
    \item \textbf{Open-source cases (5):} parser vulnerabilities in
    giflib, libxml2, tcpdump, libarchive, and Exiv2, with source code
    available and sanitizer/logic-backed validators.
\end{itemize}

The firmware cases are analyzed primarily through binary-level static
analysis, whereas the open-source cases additionally support runtime
evidence through compilation and execution.  The distinction matters
methodologically because it changes which evidence classes are available to
the attribution layer.

\subsection{Scoring}

Each case defines a scoring rubric with four check dimensions:
\begin{itemize}[leftmargin=*]
    \item \textbf{File localization:} does the agent identify the correct
    vulnerable file?
    \item \textbf{Function localization:} does the agent identify the
    correct vulnerable function?
    \item \textbf{Root cause:} does the agent describe the root cause with
    sufficient technical detail?
    \item \textbf{Evidence:} does the agent cite concrete evidence
    (function names, string references, code patterns)?
\end{itemize}
Each dimension is scored against private case metadata and aggregated onto a
0--10 scale through a weighted rubric.  The rubric is deliberately coarse:
it is intended to compare blind runs at the task level, not to certify
fine-grained semantic equivalence between two explanations.

\subsection{Evolution Service}

The evolution service consumes the feedback bundle and generates candidate
skill revisions using an LLM-based workflow engine.  Candidates are
content-deduplicated by SHA-256 hash before entering the gate queue.  The
gate runs follow-up validation and enforces the publication policy defined
in Section~\ref{sec:gate}.

\subsection{LLM Configuration}

Both the analysis agent and the evolution service use the same LLM backend
(GLM-5.2, locally deployed).  Using one model family does not eliminate
stochasticity, but it does remove cross-model capability differences as a
confound in the present study.

\section{Experimental Methodology}
\label{sec:evaluation}

\subsection{Research Questions}

\begin{description}
    \item[RQ1.] How heterogeneous are skill effects in blind vulnerability
    analysis, and can a topically relevant skill still be neutral or
    harmful?
    \item[RQ2.] Can lightweight run-level feedback serve as a useful
    triage signal for evolution, even if it does not prove causal
    contribution?
    \item[RQ3.] In a working closed loop, is candidate publication rather
    than candidate generation the main bottleneck?
\end{description}

\subsection{Benchmark Cases}

We use six firmware vulnerability cases for the A/B study
(Table~\ref{tab:cases}).  These cases cover two vulnerability types
(command injection and buffer overflow) across multiple embedded targets.
Among them, \texttt{firmware2-wireless} is a CGI command-injection task in
which the vulnerable path is exposed through a wireless-management handler,
whereas \texttt{firmware2-login} is a CGI command-injection task centered
on a login-related request path.  For readability, we refer to the three
injected skills in shorthand as \texttt{tenda-triage},
\texttt{elf-triage}, and \texttt{cgi-triage}; the corresponding full skill
artifacts remain unchanged in the repository and runtime logs.
Each case is analyzed through blind binary-level reasoning on a remote VM,
with the agent receiving neither CVE identifiers nor solution-bearing
artifact names.

\begin{table}[t]
\centering
\caption{Benchmark cases used in the A/B study.}
\label{tab:cases}
\footnotesize
\begin{tabular}{l l l c l}
\toprule
Case ID & Target & Vuln.\ type & \makecell[c]{Injected\\skill} & Evidence class \\
\midrule
f453-cmdinject & Tenda httpd & cmd injection & \makecell[l]{tenda-\\triage} & behavior \\
f453-overflow & Tenda httpd & buffer overflow & \makecell[l]{elf-\\triage} & behavior \\
f9k1122-overflow & NETGEAR httpd & buffer overflow & \makecell[l]{elf-\\triage} & behavior \\
i12-overflow & Tenda httpd & buffer overflow & \makecell[l]{elf-\\triage} & behavior \\
firmware2-wireless & Embedded CGI & cmd injection & \makecell[l]{cgi-\\triage} & behavior \\
firmware2-login & Embedded CGI & cmd injection & \makecell[l]{cgi-\\triage} & behavior \\
\bottomrule
\end{tabular}
\end{table}

\subsection{Conditions and Controls}

Each case is evaluated under two conditions:
\begin{itemize}[leftmargin=*]
    \item \textbf{No-skill (NS):} the SkillClaw server is configured to
    disable skill injection.  The agent receives only the blind prompt and
    workspace.
    \item \textbf{With-skill (WS):} the server forces injection of the
    case-specific skill listed in Table~\ref{tab:cases}.
\end{itemize}

Each condition is repeated for 3 rounds with fresh agent sessions, yielding
$6 \times 2 \times 3 = 36$ primary runs.  One additional recovery rerun was
generated operationally for a failed attempt and is archived on disk, but
it is excluded from the main A/B means reported here.  All primary runs use
the same model, inference budget, prompt format, and blind workspace
construction procedure.

\subsection{Complementary Targeted Profile Ablations}

In addition to the main six-case A/B study, we conduct a smaller
complementary ablation focused on two firmware CGI command-injection cases:
\texttt{firmware2-wireless} and \texttt{firmware2-login}.  The purpose of
this ablation is not to enlarge the benchmark, but to probe a more specific
question that emerged during development: is it enough for a skill to be
topically relevant, or does the \emph{structure} of the skill matter?

We compare four profile conditions:
\begin{itemize}[leftmargin=*]
    \item \textbf{None:} no injected firmware skill.
    \item \textbf{Wrong:} an injected firmware skill from the wrong
    procedural family (\texttt{firmware-embedded-lua-shell-extraction}).
    \item \textbf{Relevant:} the current live CGI command-injection triage
    skill, which is broader and more descriptive.
    \item \textbf{Seed:} an earlier branch-first CGI skill that prioritizes
    dispatcher variables and single-branch sink tracing.
\end{itemize}

This complementary study contributes 16 archived runs in total and is used
only for targeted interpretation in Section~\ref{sec:results_ablation}.  We
do not merge these runs into the six-case A/B means because the profile
conditions are asymmetric and were designed for mechanism diagnosis rather
than balanced cross-case comparison.

\subsection{Measurements}

We report:
\begin{itemize}[leftmargin=*]
    \item \textbf{Task quality:} 0--10 score per run, decomposed into file,
    function, root-cause, and evidence sub-scores.
    \item \textbf{Skill effect:} $\Delta = \bar{y}_{WS} - \bar{y}_{NS}$,
    where $\bar{y}$ is the mean score across rounds.
    \item \textbf{Attribution signal:} the four-level feedback status
    (supported, mismatched, contradicted, unknown) assigned after
    confirmation.
    \item \textbf{Gate outcome:} published or rejected, with rejection
    reasons.
\end{itemize}

\subsection{Statistical Analysis}

Given the small sample size ($N=3$ per condition), our analysis is
primarily descriptive.  We report individual run scores in addition to
means, and we treat non-parametric significance tests only as supporting
context rather than as the main basis for the paper's claims.

\section{Results}
\label{sec:results}

\subsection{RQ1: Skill Effects Are Unstable and Sometimes Harmful}
\label{sec:results_rq1}

Table~\ref{tab:ab_results} reports the full A/B results across all 36
primary runs.  The most important observation is heterogeneity rather than
average gain: the mean skill effect $\Delta$ ranges from $+1.67$ to
$-1.00$, and the overall average across all six cases is only $+0.28$.

\begin{table}[t]
\centering
\caption{A/B results: no-skill vs.\ with-skill across 6 cases, 3 rounds each.
$\Delta = \bar{y}_{WS} - \bar{y}_{NS}$.}
\label{tab:ab_results}
\footnotesize
\begin{tabular}{l c c c c c c l}
\toprule
Case & Skill & \makecell[c]{NS\\scores} & NS mean & \makecell[c]{WS\\scores} & WS mean & $\Delta$ & Verdict \\
\midrule
f453-cmdinject & \makecell[l]{tenda-\\triage} & 6,8,8 & 7.33 & 6,8,6 & 6.67 & $-0.67$ & Harmful \\
f453-overflow & \makecell[l]{elf-\\triage} & 6,2,6 & 4.67 & 6,2,6 & 4.67 & $\phantom{-}0.00$ & Neutral \\
f9k1122-overflow & \makecell[l]{elf-\\triage} & 2,6,6 & 4.67 & 6,6,6 & 6.00 & $+1.33$ & Helpful \\
i12-overflow & \makecell[l]{elf-\\triage} & 3,8,3 & 4.67 & 3,8,8 & 6.33 & $+1.67$ & Helpful \\
\makecell[l]{firmware2-\\wireless} & \makecell[l]{cgi-\\triage} & 5,5,5 & 5.00 & 5,5,2 & 4.00 & $-1.00$ & Harmful \\
\makecell[l]{firmware2-\\login} & \makecell[l]{cgi-\\triage} & 5,4,5 & 4.67 & 5,5,5 & 5.00 & $+0.33$ & Marginal \\
\midrule
\multicolumn{5}{l}{Average across all cases} &  & $+0.28$ &  \\
\bottomrule
\end{tabular}
\end{table}

\textbf{Finding 1: Skill effects are highly case-dependent.}
The same skill (\texttt{elf-cwe120-firmware-triage}) produces $+1.33$ on
f9k1122-overflow, $+1.67$ on i12-overflow, but $0.00$ on f453-overflow.
This result cautions against any blanket statement such as ``domain skill X
improves binary analysis.''  A skill may encode generally sensible tactics
yet still fail to help a particular case.

\textbf{Finding 2: Topically relevant skills can still degrade
performance.}
The firmware2-wireless case injected \texttt{cgi-triage}, a
skill that is topically relevant to the task (CGI command injection in
embedded firmware).  Yet the with-skill condition produced a lower mean
score ($4.00$ vs.\ $5.00$), with one round dropping to $2$.  This mirrors
recent evidence that many skill-induced failures are caused not by
obviously irrelevant instructions, but by superficially plausible skills
that distort the agent's reasoning trajectory~\cite{skillharmful}.

\textbf{Finding 3: run-to-run model variance is large enough to obscure
moderate skill effects.}
The no-skill condition for i12-overflow produces scores of $3, 8, 3$---a
range of 5 points across three identical-configuration runs.  Even though
the with-skill mean is $1.67$ points higher, this difference is not
statistically significant at $N=3$ (Wilcoxon $p > 0.1$).  This demonstrates
that LLM output variance is a major confound for skill attribution: even
without any skill, the model's performance fluctuates enough to mask
moderate skill effects.

\textbf{Finding 4: some injected skills appear observationally redundant.}
The f453-overflow case shows identical scores in both conditions
($6, 2, 6$), indicating that the injected skill \texttt{elf-cwe120-firmware-triage}
had no observable effect on the agent's behavior.  This may mean the agent
ignored the skill, or that the skill's content was redundant with the
model's existing knowledge.  Distinguishing these cases requires
trajectory-level analysis (Section~\ref{sec:discussion}).

\subsection{Complementary Result: Targeted CGI Profile Ablations}
\label{sec:results_ablation}

The six-case A/B study tells us whether a chosen skill helps on average
relative to a no-skill baseline.  It does not reveal whether a skill's
\emph{topic match} is enough, whether a narrower seed version can be more
effective than a broader live version, or whether a wrong-domain firmware
skill can still produce plausible outputs.  We therefore inspect the 16-run
targeted CGI ablations introduced in
Section~\ref{sec:evaluation}.

\begin{table}[t]
\centering
\caption{Complementary profile ablations on two firmware CGI cases.  Values
are mean scores over archived runs under each profile condition.}
\label{tab:profile_ablation}
\footnotesize
\begin{tabular}{lcccc}
\toprule
Case & None & Wrong & Relevant & Seed \\
\midrule
\texttt{firmware2-wireless} & 2.67 & 3.00 & 3.00 & \textbf{4.00} \\
\texttt{firmware2-login} & 4.33 & 4.50 & \textbf{5.00} & \textbf{5.00} \\
\bottomrule
\end{tabular}
\end{table}

Two observations matter.

\textbf{First, topical relevance alone is insufficient.}  On
\texttt{firmware2-wireless}, the broader live \texttt{relevant} profile and
the wrong-domain firmware profile both score $3.00$, only slightly above
the no-skill mean of $2.67$.  A skill can therefore sound domain-relevant
without actually providing the procedural bias needed by the case.

\textbf{Second, skill structure matters.}  The branch-first
\texttt{seed} profile reaches the highest mean on
\texttt{firmware2-wireless} ($4.00$) and ties the best mean on
\texttt{firmware2-login} ($5.00$).  The difference between
\texttt{seed} and \texttt{relevant} is not topical; both target CGI
command-injection analysis.  The difference is procedural emphasis.  The
\texttt{seed} version first anchors the dispatcher variable (e.g.,
\texttt{page}/\texttt{action}/\texttt{mode}) and then traces one concrete
branch to the sink, whereas the broader live profile allocates more prompt
budget to generic CGI entrypoint enumeration and request structure.

\textbf{Interpretation.}  These ablations do not yet prove strict skill
necessity: on \texttt{firmware2-login}, both \texttt{relevant} and
\texttt{seed} reach the same mean, and even the \texttt{wrong} profile
remains close to the no-skill baseline.  What they do show is that
``correct topic'' and ``helpful skill'' are not the same notion.  For this
kind of firmware CGI task, the current evidence suggests that branch-first
procedural guidance is more stable than broad descriptive guidance, but the
effect is still case-dependent and not yet strong enough to establish an
all-or-nothing necessity threshold.

\subsection{RQ2: Evidence-Conditioned Feedback vs.\ A/B Attribution}
\label{sec:results_rq2}

We next compare the four-level feedback signal with the A/B $\Delta$
measure.  The purpose of this comparison is not to prove that the feedback
signal is causal; rather, it is to test whether the signal at least points
in the correct directional regime.  For the six firmware cases in this
study:

\begin{itemize}[leftmargin=*]
    \item \textbf{i12-overflow ($\Delta = +1.67$):} The skill
    \texttt{elf-triage} describes ELF binary analysis
    steps that align with the agent's trajectory (function localization
    via PLT/GOT, sink identification via \texttt{system}/\texttt{sprintf}).
    Feedback: \emph{supported}.  This is consistent with the positive
    $\Delta$.
    \item \textbf{firmware2-wireless ($\Delta = -1.00$):} The skill
    \texttt{cgi-triage} describes CGI handler
    analysis, which is relevant to the task.  However, the with-skill run
    in round 3 dropped to score 2 (file miss), suggesting the skill may
    have directed the agent to the wrong binary or function.  Feedback:
    \emph{contradicted}.  Consistent with the negative $\Delta$.
    \item \textbf{f453-overflow ($\Delta = 0.00$):} The skill was injected
    but had no effect on scores.  Feedback: \emph{mismatched} (the skill
    is topically related to ELF analysis, but the trajectory shows no
    evidence of the agent following the skill's specific steps).
    Consistent with $\Delta \approx 0$.
\end{itemize}

Taken together, these cases suggest that the four-level signal is
\emph{directionally plausible}: supported labels tend to coincide with
positive A/B effects, contradicted labels with negative effects, and
mismatched labels with near-zero effects.  At the same time, the signal is
best viewed as a triage mechanism rather than a proof device.  Its value is
that it can prioritize which skills deserve ablation, rejection, or further
inspection without paying the full cost of paired reruns for every case.

\subsection{RQ3: Replay-Only Gate Has a Structural Limitation}
\label{sec:results_gate}

We analyze 94 gate decisions collected across all experiments
(Table~\ref{tab:gate_stats}).

\begin{table}[t]
\centering
\caption{Gate decision summary across all experiments.}
\label{tab:gate_stats}
\small
\begin{tabular}{lr}
\toprule
Metric & Count \\
\midrule
Total gate decisions & 94 \\
Published & 12 \\
Rejected & 82 \\
Published under earlier score-only gate logging & 10 \\
Published under baseline-aware non-inferiority logging & 2 \\
\bottomrule
\end{tabular}
\end{table}

The archived decision history spans two gate regimes.  Earlier decisions
record only the candidate-side validation score and final decision.  Later
decisions additionally record a baseline mean, gate mode, and effective
threshold.  Among the latter, the two published examples were accepted
under a non-inferiority criterion, with $(0.7, 0.4)$ and $(0.6, 0.6)$ for
(candidate mean, baseline mean), respectively.

\textbf{Root cause analysis.}
The key issue is not simply that many candidates are poor.  Rather, the
gate often lacks a defensible basis for deciding whether a candidate skill
was \emph{necessary}.  In a replay-only setting, a candidate can reproduce
the same validated answer as the baseline while providing no clear evidence
that it changed the outcome.  Under a strict-improvement interpretation,
such a candidate is rejected as redundant; under a non-inferiority
interpretation, the same candidate may be retained as harmless but
potentially useful.  The decision therefore depends as much on the gate's
philosophy as on the candidate's intrinsic quality.

This is the structural limitation formalized in
Section~\ref{sec:publication_problem}.  The gate's
$\mathrm{FollowupPass}$ function computes $q_{\mathrm{task}}(y_{\Delta s})$
but only imperfectly captures $q_{\mathrm{task}}(y_0)$.  When the model can
already solve the task without the skill, a replay-only gate struggles to
separate three cases: genuine improvement, harmless redundancy, and
spurious success.  This ambiguity is exactly why publication remains the
hardest part of the loop.

\textbf{Interpretation.}
Recent work on skill misevolution~\cite{misevolution} proposes SafeEvolve, a wrapper that repairs unsafe
content and governs subsequent reuse.  Our study highlights a closely
related but distinct problem: even when content appears safe, the system
may still be unable to decide whether the skill deserves to persist.  In
other words, safety and usefulness are separable questions, and a realistic
evolution loop must reason about both.

\subsection{Summary of Findings}

\begin{table}[t]
\centering
\caption{Summary of findings mapped to hypotheses.}
\label{tab:findings}
\small
\begin{tabularx}{\textwidth}{lXl}
\toprule
Hypothesis & Finding & Status \\
\midrule
H1: Skill instability & $\Delta$ ranges from $-1.00$ to $+1.67$; 2/6 cases show net harm & Supported \\
H2: Evidence approximation & Four-level feedback is directionally consistent with A/B $\Delta$ & Partially supported \\
H3: Gate limitation & 82/94 candidates rejected; baseline-aware decisions remain difficult to interpret without counterfactual evidence & Supported \\
\bottomrule
\end{tabularx}
\end{table}

\section{Discussion}
\label{sec:discussion}

\subsection{The Evidence--Attribution Gap}

Our results reveal what we call the \emph{evidence--attribution gap}:
executable domain evidence can verify \emph{what} the agent found (the
vulnerability is real, the localization is correct), but it cannot
directly verify \emph{why} the agent found it (was the skill causally
responsible?).

This gap has two dimensions:
\begin{enumerate}[leftmargin=*]
    \item \textbf{Positive evidence is not causal evidence.} A run that
    passes source-backed validation and has a ``supported'' feedback status
    may have succeeded because of the skill, because of the model's prior
    knowledge, or because of both.  The four-level feedback reduces this
    ambiguity by checking skill--trajectory alignment, but it cannot
    eliminate it.
    \item \textbf{Absence of evidence is not evidence of absence.} A
    ``mismatched'' skill may have subtly influenced the agent's reasoning
    in ways that the alignment check cannot detect.  For example, a skill
    might prime the agent to look at certain binary sections, indirectly
    leading to the correct answer even if the skill's specific steps were
    not followed.
\end{enumerate}

Bridging this gap requires either counterfactual ablation (expensive but
precise) or a probabilistic attribution model that estimates $P(\text{success}
\mid \text{skill}) - P(\text{success} \mid \text{no skill})$ from observed
correlations (cheaper but requires more data).  We leave the latter as
future work.

\subsection{Implications for Skill Library Design}

Our finding that the same skill can be helpful, neutral, or harmful
depending on the case has a direct implication: \emph{skill libraries need
case-sensitive selection and case-sensitive evaluation.}  A globally
plausible skill may still degrade a specific analysis if its procedural
biases do not fit the target.  This is consistent with
prior observations that skill-induced failures are strongly
case-specific~\cite{skillharmful}.

This suggests that retrieval mechanisms should consider not only textual
skill--task similarity, but also uncertainty about the model's unaided
capability.  If the baseline is already strong, injection may add noise,
encourage over-commitment to the wrong pattern, or simply consume prompt
budget without benefit.

\subsection{Toward a Counterfactual Gate}
\label{sec:counterfactual_gate}

The structural limitation of the replay-only gate
(Section~\ref{sec:results_gate}) suggests two directions for improvement:

\textbf{Scheme A: Counterfactual ablation gate.} For each candidate, run
the task both with and without the candidate skill, and require strict
improvement ($\bar{y}_{WS} > \bar{y}_{NS}$).  This directly measures
$\Delta(\Delta s)$ but doubles the validation cost and requires
statistical power (multiple rounds).

\textbf{Scheme B: Non-inferiority gate.} Instead of requiring strict
improvement, require that the candidate is \emph{not worse} than the
baseline by more than a margin $\epsilon$: $\bar{y}_{WS} \geq
\bar{y}_{NS} - \epsilon$.  This relaxes the gate to allow redundant (but
not harmful) candidates, at the cost of potentially publishing skills that
do not help.

A combined approach could use Scheme B for initial publication (allowing
the skill to enter the live library) and Scheme A for retention decisions
(removing skills that do not demonstrate improvement over a longer
horizon).  More broadly, this suggests that publication and retention may
need to be decoupled: the system can admit a candidate as provisionally
safe while still demanding stronger evidence before treating it as
established expertise.

\subsection{Comparison with Concurrent Work}

Table~\ref{tab:comparison} compares our approach with concurrent work on
agent skills.

\begin{table}[t]
\centering
\caption{Comparison with concurrent work on agent skills.}
\label{tab:comparison}
\footnotesize
\begin{tabular}{l c c c c c}
\toprule
& \multicolumn{2}{c}{Attribution} & \multicolumn{2}{c}{Publication} & \\
\cmidrule(lr){2-3} \cmidrule(lr){4-5}
Approach & Method & Cost & Method & Safety & Domain \\
\midrule
Ours & 4-level evidence & 1 run & Replay gate & Conservative & Vuln. analysis \\
SkillShapley~\cite{skillshapley} & Shapley value & $O(2^n)$ & N/A & N/A & General \\
Skill Harmful~\cite{skillharmful} & Differential & 2 runs & N/A & N/A & General \\
Misevolution~\cite{misevolution} & N/A & N/A & SafeEvolve & Repair-based & General \\
\bottomrule
\end{tabular}
\end{table}

Our approach is distinctive in its use of executable domain evidence as a
post-run attribution signal.  This gives it a cost advantage over explicit
ablation, but also a narrower claim: we obtain an auditable approximation,
not a full causal decomposition.

\subsection{Limitations of Static Analysis}

The firmware cases in our A/B study are analyzed through binary-level
static reasoning without a live runtime environment.  Consequently, the
available evidence is weaker than sanitizer-backed confirmation in the
open-source cases.  This likely lowers the precision of the feedback signal
and should be interpreted as a limitation of the present study, not as a
property of the framework in general.

\section{Related Work}
\label{sec:related}

\subsection{Skill Induction, Libraries, and Evolution}

Voyager~\cite{voyager}, ExpeL~\cite{expel}, Reflexion~\cite{reflexion},
and MemGPT~\cite{memgpt} establish the broader agent-learning context in
which trajectories, memories, and reusable instructions are accumulated
across tasks.  SkillClaw~\cite{skillclaw} makes this logic explicit at the
level of a shared natural-language skill library, while
Ctx2Skill~\cite{ctx2skill} and OPID~\cite{opid} study, respectively,
context-to-skill induction and trajectory-derived skill distillation.
These works show that skill formation and reuse are already central to LLM
agent research.  Our paper does not compete by proposing yet another
generic skill generator.  Instead, it focuses on a narrower systems
question: in a vulnerability-analysis deployment, how should a completed
run be transformed---or refused transformation---into a live reusable
skill?

\subsection{Environment- and Verifier-Grounded Skill Refinement}

Recent work increasingly grounds skill refinement in environmental or
verifier feedback rather than in free-form self-reflection alone.
SPARK~\cite{spark} emphasizes \emph{evidence over plans}, preferring
environment-grounded revisions to purely narrative improvement.
Trace2Skill~\cite{trace2skill} distills trajectory-local lessons into
longer-term skills using verifier feedback, and
VeriSkill~\cite{veriskill} studies self-evolution specifically for program
verification skills.  SkillEvolBench~\cite{skillevolbench} further turns
episodic-to-procedural evolution into an explicit benchmark problem.

Our work is aligned with this trend, but differs in scope and target.  We
do not study general verifier-backed skill improvement in the abstract.
We study the security-specific transition from \emph{blind remote
vulnerability-analysis runs} to \emph{candidate skills that may alter a
live library}.  In this setting, the central issue is not only whether
verifier signals exist, but how they participate in archival, candidate
staging, and publication control.

\subsection{Skill Failure, Attribution, and Governance}

Skill-induced failures~\cite{skillharmful} demonstrate that apparently
relevant skills can actively harm agent performance by steering execution
onto the wrong procedure.  Skill misevolution~\cite{misevolution} shows
that self-improving loops can preserve unsafe behavior as persistent
skills.  Adversarial skill publishing~\cite{detourhijack} further exposes
how untrusted skills can manipulate downstream behavior.

These works motivate our emphasis on post-run governance.  In our setting,
retrieval and generation are only the front half of the problem.  The back
half is deciding whether a candidate skill deserves admission into the live
library.  Our gate analysis complements prior work by showing that even
when candidate generation is available, publication remains difficult
because replay-only evidence is often insufficient to separate improvement
from redundancy.

\subsection{Skill Attribution and Valuation}

SkillShapley~\cite{skillshapley} formulates skill-step attribution through
Shapley-value estimation~\cite{shapley}, yielding a principled but costly
counterfactual approach.  Differential no-skill / with-skill
comparison~\cite{skillharmful} provides a cheaper alternative at the run
level.  Our feedback interface belongs to this family of post-run
assessment methods, but with a deliberately narrower claim: it is a
triage-oriented operational signal, not a full causal estimator.  Its
purpose is to support candidate generation and publication decisions in a
security workflow where full ablations are expensive and often delayed.

\subsection{Security Agents, Playbooks, and Benchmarks}

On the security side, PentestGPT~\cite{pentestgpt} and subsequent studies
on LLM-based vulnerability analysis~\cite{gpt4vuln,aslr} focus on what
models can discover or explain.  More recent work such as
EvoHunt~\cite{evohunt} begins to evolve reusable playbooks for agentic
security auditing, while SEC-bench Pro~\cite{secbenchpro} raises the bar on
realistic long-horizon software-security evaluation.  Our work is adjacent
to both directions.  Like EvoHunt, we care about reusable procedural
knowledge in security tasks; like SEC-bench Pro, we care about realistic
task construction and leakage control.  The difference is that we center
the run-to-library transition itself: blind execution, external
confirmation, candidate staging, and live-skill admission are the primary
objects of study.

\section{Threats to Validity}
\label{sec:threats}

\paragraph{Sample size.}
The A/B study uses $N=3$ rounds per condition, which provides limited
statistical power.  The Wilcoxon signed-rank test does not reach
significance for any individual case.  While the directional consistency
of the four-level feedback with A/B $\Delta$ is suggestive, a larger study
($N \geq 10$) is needed for statistical confirmation.

\paragraph{Model variance.}
LLM outputs are stochastic.  The i12-overflow no-skill condition produces
scores of $3, 8, 3$, demonstrating a 5-point range under identical
configuration.  This variance is a confound for attribution: moderate skill
effects ($\Delta \approx 1$--$2$) may be masked by model noise.  Future
work should either increase $N$ to average out variance or use
temperature-0 decoding for reproducibility (at the cost of ecological
validity).

\paragraph{Single model.}
All experiments use one model (GLM-5.2).  Skill effects may differ across
models, as different models have different prior knowledge and instruction
following capabilities.  Cross-model validation is left as future work.

\paragraph{Static analysis limitation.}
The firmware cases in the main A/B study are analyzed through binary-level
static reasoning without a live runtime environment.  Consequently, the
available evidence is weaker than sanitizer-backed confirmation in the
open-source cases.  This likely lowers the precision of the feedback signal
and should be interpreted as a limitation of the present study, not as a
property of the framework in general.

\paragraph{Attribution precision.}
The four-level feedback is a heuristic approximation, not a causal
estimator.  ``Supported'' does not prove causality, and ``mismatched''
does not prove the absence of contribution.  A rigorous validation would
require comparing the four-level signal against ground-truth causal
effects (from counterfactual ablation) across many cases.

\paragraph{Gate evaluation.}
The 94 gate decisions were produced under multiple development-time gate
regimes.  Earlier decisions preserve only candidate-side validation
outcomes, whereas later decisions include baseline-aware metadata and
non-inferiority parameters.  This historical mixture is useful for
retrospective diagnosis but imperfect for a clean scientific comparison.  A
future study should evaluate gate variants from fixed repository snapshots
under a single benchmark protocol.

\paragraph{Answer leakage.}
Known-CVE benchmarks can leak through identifiers, artifact names,
comments, or pre-existing PoCs.  While our blind workspaces are generated
automatically and audited to exclude solution-bearing material, we cannot
guarantee zero leakage, especially for well-known CVEs.

\paragraph{Oracle validity.}
Case validators encode trusted expectations and may themselves be
incomplete.  A validator that passes does not guarantee the vulnerability
is fully understood; a validator that fails does not guarantee the analysis
is wrong.  We report validator provenance and evidence class to allow
readers to assess the strength of each confirmation.

\section{Conclusion}
\label{sec:conclusion}

We studied the transition from blind vulnerability-analysis runs to live
reusable skills.  The main claim of the paper is not that autonomous skill
improvement is already solved.  Rather, we show that this transition can be
made explicit enough to study: blind remote execution, local external
confirmation, run archival, candidate generation, and publication control
can be connected into a single operational loop.

Empirically, the picture is mixed in a useful way.  Across 36 remote blind
A/B runs, skill effects are strongly heterogeneous: the same broad class of
skill can help one case, do nothing on another, and harm a third.
Complementary CGI profile ablations further suggest that procedural
structure matters at least as much as topical relevance.  A branch-first
skill can outperform a broader descriptive variant on one firmware case
while tying it on another, yet the current evidence still falls short of a
strict necessity claim.

We also show that publication control is the real bottleneck of the
current loop.  The system can already generate candidate skills and replay
them against archived cases, but replay-only evidence is often not strong
enough to justify admission into the live library.  In practice, the hard
problem is not writing another candidate skill.  The hard problem is
deciding whether a validated run contains knowledge that should persist.

This observation is relevant beyond vulnerability analysis.  If the
run-to-library transition is difficult even in a domain with executable
external confirmation, it is likely harder in agent settings where such
confirmation is weak or absent.  Our framework suggests a pragmatic stance:
external evidence can discipline skill evolution, but it does not remove
the need for explicit publication and retention policies.

Future work should proceed along four directions: (1) scale the blind study
to larger case sets and more rounds; (2) implement counterfactual or
retention-stage gates that better separate improvement from redundancy; (3)
incorporate stronger runtime-backed evidence on open-source cases; and (4)
test whether the same lifecycle design transfers to broader security-agent
settings such as exploit validation or multi-step auditing playbooks.

\section*{Acknowledgments}
Acknowledgments go here.

\bibliographystyle{elsarticle-num}
\bibliography{references}

\end{document}
